\chapter{Stato dell'arte e Background}
\label{cha:intro}

L'ultimo decennio ha segnato un punto di svolta fondamentale nella storia
dell'informatica, caratterizzato da una transizione senza precedenti
verso sistemi basati sull'Intelligenza Artificiale (AI). Quella che
inizialmente era considerata una disciplina di nicchia è evoluta in un
pilastro tecnologico pervasivo, capace di influenzare settori che
spaziano dalla medicina alla finanza, fino alla gestione dei dati
strutturati.

\section{Recenti scoperte ed invenzioni}\label{recenti-scoperte-ed-invenzioni}

Nell'ultimo periodo abbiamo assistito alla nascita dei Large Language
Models (LLM), come la famiglia GPT di OpenAI, Claude di Anthropic e
modelli open-source come Llama di Meta. Questi sistemi non si limitano
più a una mera elaborazione statistica del linguaggio, ma dimostrano
capacità emergenti di ragionamento logico, sintesi e, in particolare,
generazione di codice. È proprio in questo alveo di innovazione che si
inserisce la capacità di tradurre il linguaggio naturale in linguaggi di
programmazione o di interrogazione, aprendo la strada a interfacce
uomo-macchina sempre più intuitive e meno dipendenti dalla conoscenza
tecnica sintattica.

L'avvento degli LLM ha tuttavia evidenziato alcuni limiti intrinseci
legati alla natura statica del loro addestramento e alla tendenza a
generare informazioni errate ma sintatticamente plausibili (fenomeno
noto come allucinazione). Per superare tali vincoli e connettere
i modelli linguistici a basi di conoscenza dinamiche, la ricerca in oggetto si è
focalizzata sullo sviluppo di architetture capaci di operare nello
spazio della semantica computazionale.

Il cardine di questa evoluzione è rappresentato dagli embeddings, ovvero
trasformazioni matematiche che convertono dati discreti, come parole,
frasi o immagini, in vettori numerici continui all'interno di uno spazio
a dimensionalità ridotta \cite{geeksforgeeks2024embeddings}. L'obiettivo fondamentale di queste
rappresentazioni è codificare le informazioni semantiche in modo che
elementi concettualmente simili risultino matematicamente vicini
all'interno dello spazio vettoriale, calcolandone la distanza geometrica
\cite{wang-etal-2024-improving-text}. A differenza dei metodi tradizionali basati sulla
corrispondenza testuale esatta o sulle parole chiave, gli embeddings
catturano relazioni complesse e sfumature del linguaggio, permettendo al
sistema di comprendere il contesto e il significato profondo dell'input
\cite{thinkia2025vectordatabases}.

Su questa infrastruttura matematica si innesta il paradigma della
Retrieval-Augmented Generation (RAG), una metodologia che potenzia i
modelli linguistici integrando la loro memoria nativa (memoria
parametrica) con un modulo di recupero della conoscenza esterna
(memoria non parametrica) \cite{lewis2021retrievalaugmentedgenerationknowledgeintensivenlp}. Il processo si articola
attraverso la sinergia di due componenti principali \cite{gupta2024comprehensivesurveyretrievalaugmentedgeneration}:

\begin{enumerate}
\def\labelenumi{\arabic{enumi}.}
\item
  Il modulo di recupero (Retriever): I documenti della base di
  conoscenza esterna vengono frammentati, convertiti in embeddings e
  indicizzati. All'atto dell'interrogazione, anche la query dell'utente
  viene trasformata in un vettore, consentendo al sistema di
  confrontarla con i documenti indicizzati e identificare le unità
  informative semanticamente più pertinenti.
\item
  Il modulo di generazione (Generator): Il contesto recuperato viene
  inserito nel prompt originale e fornito all'LLM, il quale utilizza
  queste informazioni aggiornate per produrre una risposta accurata e
  fattualmente fondata.
\end{enumerate}

Sotto il profilo metodologico, l'approccio RAG offre tre vantaggi
cruciali nell'ingegneria dei sistemi basati su AI. In primo luogo,
l'ancoraggio a documenti reali mitiga drasticamente il rischio di
allucinazioni \cite{lewis2021retrievalaugmentedgenerationknowledgeintensivenlp, gupta2024comprehensivesurveyretrievalaugmentedgeneration, enwiki:1360040196}. In secondo luogo, permette un
aggiornamento dinamico della conoscenza, poiché la memoria del sistema
può essere espansa o revisionata modificando semplicemente l'indice dei
vettori, senza l'oneroso bisogno di riaddestrare il modello \cite{gupta2024comprehensivesurveyretrievalaugmentedgeneration,
lewis2021retrievalaugmentedgenerationknowledgeintensivenlp}. Infine, garantisce un'elevata efficienza computazionale,
ottimizzando l'uso della finestra di contesto del modello generativo
attraverso la selezione esclusiva dei sottoinsiemi di dati rilevanti
\cite{glass2025extractiveschemalinkingtexttosql}.

Da un punto di vista infrastrutturale, la gestione e l'interrogazione
efficiente di milioni di vettori ad alta dimensionalità richiede
tecnologie dedicate: i vector store (o database vettoriali) \cite{thinkia2025vectordatabases,
enwiki:1360040196}. Questa tecnologia di archiviazione alimenta la ricerca
semantica implementando algoritmi avanzati per la ricerca dei vicini più
prossimi (come l' Approximate Nearest Neighbor), capaci di
restituire i vettori più simili in tempi estremamente ridotti \cite{khalusova2024understanding,
enwiki:1360040196}. All'interno dei sistemi RAG, il vector store funge da vero e
proprio ponte logistico tra l'LLM e i repository di dati; esso memorizza
i vettori e gestisce i relativi metadati, facilitando il filtraggio
granulare delle informazioni e migliorando la trasparenza del sistema
grazie alla tracciabilità e alla citazione delle fonti originali
\cite{thinkia2025vectordatabases, enwiki:1360040196}.

\section{Metodologia Text-to-SQL e Integrazione del
RAG}\label{metodologia-text-to-sql-e-integrazione-del-rag}

La traduzione Text-to-SQL rappresenta un compito cardine del
Natural Language Processing (NLP), il cui fine è la conversione
automatica di quesiti in linguaggio naturale (NL) in interrogazioni
SQL eseguibili da un sistema
di gestione di database relazionali (RDBMS) \cite{hong2025nextgenerationdatabaseinterfacessurvey, kanburoglu2024text}. Questo
processo funge da interfaccia intuitiva verso i dati, consentendo l'estrazione di
informazioni anche a utenti privi di competenze tecniche \cite{lei-etal-2020-examining, xia2024r3thissqlme,
hong2025nextgenerationdatabaseinterfacessurvey}.

Sotto il profilo computazionale, il problema è modellato come una sfida
Sequence-to-Sequence (\(Seq2Seq\)), in cui l'input unisce il
quesito dell'utente e lo schema relazionale del database per generare la
stringa SQL corrispondente \cite{xia2024r3thissqlme, kanburoglu2024text}.

Nonostante i progressi dei Large Language Models (LLM), la
traduzione accurata sconta criticità legate all'ambiguità semantica e
alla scalabilità degli schemi. Nei database reali, la presenza di
centinaia di tabelle rende proibitivo l'inserimento dell'intero catalogo
nel prompt a causa dei limiti della finestra di contesto e dei relativi costi computazionali \cite{hong2025nextgenerationdatabaseinterfacessurvey}. In questo
scenario, l'integrazione di tecniche di Retrieval-Augmented
Generation (RAG) costituisce una strategia determinante.

L'applicazione del paradigma RAG al Text-to-SQL si realizza
principalmente attraverso lo Schema Linking, volto ad allineare i
termini della domanda con le entità specifiche del database \cite{lei-etal-2020-examining,
glass2025extractiveschemalinkingtexttosql, nahid2026rethinkingschemalinkingcontextaware}. Invece di trasmettere l'intero catalogo, un modulo di
recupero (retriever) seleziona esclusivamente il sottoinsieme di
tabelle e colonne pertinente al quesito \cite{glass2025extractiveschemalinkingtexttosql, peng2025xsqlexpertschemalinking}. Tale approccio
isola uno ``schema focalizzato'',
riducendo il rumore informativo e prevenendo i fenomeni di allucinazione
del modello (ovvero il riferimento a identificatori inesistenti)
\cite{glass2025extractiveschemalinkingtexttosql, nahid2026rethinkingschemalinkingcontextaware}.

I principali vantaggi tecnici e metodologici di questa integrazione
includono:

\begin{itemize}
\item
  \textbf{efficienza dei token e contrazione dei costi:} il filtraggio
  preventivo dello schema riduce il volume di token inviati all'LLM,
  abbattendo la latenza e ottimizzando l'impiego di modelli
  open-weight minori; \cite{glass2025extractiveschemalinkingtexttosql, peng2025xsqlexpertschemalinking, nahid2026rethinkingschemalinkingcontextaware}
\item
  \textbf{precisione nello Schema Linking:} l'uso di embedding
  vettoriali e ricerca semantica cattura relazioni latenti tra i
  concetti espressi dall'utente e la nomenclatura tecnica, superando i
  limiti della corrispondenza lessicale; \cite{thinkia2025vectordatabases, geeksforgeeks2024embeddings}
\item
  \textbf{robustezza Cross-Domain:} il framework facilita la
  generalizzazione in modalità zero-shot su domini mai visti in
  fase di addestramento, estraendo dinamicamente metadati aggiornati
  senza richiedere il riaddestramento della rete; \cite{wang-etal-2024-improving-text, hong2025nextgenerationdatabaseinterfacessurvey}
\item
  \textbf{strategie di recupero avanzate:} approcci combinati di tipo
  table-first (identificazione delle tabelle e poi delle colonne)
  e column-first (selezione degli attributi per risalire alle
  tabelle) garantiscono un'elevata recall statistica, preservando
  gli elementi essenziali per la query. \cite{nahid2026rethinkingschemalinkingcontextaware}
\end{itemize}

L'integrazione con le tecniche RAG trasforma così il Text-to-SQL da una
generazione stocastica ``alla cieca'' a un processo deterministico
guidato dai dati.

Oltre all'ottimizzazione del contesto, la ricerca recente ha introdotto
guide metodologiche esplicite per strutturare i processi inferenziali
del modello, incrementando l'accuratezza rispetto al prompting
diretto.

La tecnica cardine è il Chain-of-Thought (CoT): l'LLM viene
istruito a esplicitare il proprio ragionamento step-by-step,
scomponendo il problema in passaggi logici intermedi (individuazione
delle tabelle, applicazione dei filtri e risoluzione delle giunzioni).
\cite{xia2025chainofthoughtsurveychainofxparadigms} Questa scomposizione riduce gli errori logico-strutturali e
ottimizza la gestione delle clausole di \texttt{JOIN} complesse.

Infine, il concetto di Self-Correction istituisce un ciclo
iterativo a ciclo chiuso in cui il sistema agisce come revisore del
proprio output \cite{gupta2024comprehensivesurveyretrievalaugmentedgeneration, rahmani2025selfcorrectinglargelanguagemodels}. Qualora l'esecuzione sul DBMS
sollevi un'eccezione sintattica o di runtime, il modello riceve in input
il messaggio d'errore nativo e utilizza questo feedback informativo per
correggere la query e riavviare l'esecuzione \cite{gupta2024comprehensivesurveyretrievalaugmentedgeneration}. Tale
meccanismo consente di sanare query inizialmente non conformi,
innalzando l'accuratezza esecutiva rispetto
ai sistemi tradizionali a tentativo unico.

\section{Approfondimento sui tool presenti in
commercio}\label{approfondimento-sui-tool-presenti-in-commercio}

È interessante notare come il mercato si sia diviso in due segmenti
chiari: tool orientati alla Business Intelligence (BI) per data analyst
e framework orientati alla produttività dello sviluppatore.

L'attuale ecosistema Text-to-SQL offre una vasta gamma di interfacce web
e agenti AI che mirano a democratizzare l'accesso ai dati. L'analisi che
segue classifica i tool in base alla loro capacità di integrazione e
alla raffinatezza del loro approccio al problema del contesto.

Tra le soluzioni più snelle spicca SQLAI.ai, che fa della flessibilità
nella gestione dello schema il suo tratto distintivo; permettendo sia
l'input manuale che la connessione diretta, si pone come un supporto
versatile per il programmatore, specialmente in contesti dove la
sicurezza del dato è prioritaria \cite{sqlai}.

Diversa è invece la strategia adottata da Chat2DB, che punta
sull'eterogeneità del supporto: la sua architettura permette di
interfacciare il sistema con una vasta gamma di modelli linguistici, da
Claude a GPT, fino a modelli proprietari via API, rendendolo uno
strumento agnostico rispetto al fornitore di IA e adatto ad ambienti
enterprise complessi \cite{chat2db}.

Sul fronte della personalizzazione e del recupero dell'informazione,
Vanna.ai rappresenta lo stato dell'arte dei framework basati su RAG. A
differenza di una semplice web-app, Vanna opera come una libreria Python
che sfrutta una memoria semantica per apprendere dai metadati e dalle
esecuzioni passate, fornendo all'utente una visione trasparente del
processo logico \cite{vanna_ai}. Infine, l'evoluzione verso agenti autonomi è incarnata
da AskYourDatabase \cite{askyourdatabase}. Attraverso l'introduzione della cosiddetta `Think
Mode', il sistema esplicita il ragionamento Chain-of-Thought del
modello, posizionandosi non solo come generatore SQL, ma come una vera e
propria piattaforma di Business Intelligence conversazionale pronta
all'integrazione via API.

Questi strumenti mirano a eliminare completamente la barriera del
linguaggio SQL, offrendo insight diretti e visualizzazioni. Sia
Text2SQL.ai che AI2sql si focalizzano sulla semplicità d'uso per i non
tecnici \cite{text2sql,ai2sql}. Mentre Text2SQL.ai eccelle nella spiegazione testuale del
codice generato, AI2sql offre strumenti unici per la creazione di schemi
da zero (tramite file SQL o descrizioni naturali), risultando utile
anche in fase di progettazione del database. BlazeSQL, si posiziona come
un assistente di alto livello per team aziendali \cite{blazesql}. Include funzionalità
di addestramento specifico sullo schema utente tramite richieste di
prova. Questo suggerisce l'uso di tecniche di Few-Shot Prompting per
migliorare la precisione nel tempo in base al dominio specifico
dell'azienda. Formula Bot e Sequel estendono il concetto di Text-to-SQL
all'analisi documentale \cite{formulabot,sequel}. Formula Bot, in particolare, permette di unire
dati provenienti da database SQL con file CSV e PDF, offrendo una
visione olistica dei dati aziendali e la possibilità di scaricare i
record analizzati. SQLchat adotta un approccio ``Bring Your Own Key'': 
non fornisce un modello interno, ma richiede le chiavi API
dell'utente (es. OpenAI) \cite{sqlchat}. Questo garantisce all'utente il controllo
totale sui costi e sulle versioni del modello utilizzato, pur mantenendo
un'interfaccia di chat intuitiva connessa al database.

Dall'analisi emerge che il fattore discriminante tra un tool
``giocattolo'' e uno professionale è la gestione del contesto (Schema
Linking). Mentre i tool più semplici chiedono all'utente di specificare
le tabelle (approccio manuale), i sistemi più evoluti come Vanna,
AskYourDatabase e BlazeSQL implementano internamente architetture RAG
per selezionare autonomamente le informazioni pertinenti, riflettendo le
metodologie di ricerca descritte nei capitoli precedenti.

L'analisi dello stato dell'arte evidenzia come il passaggio dal
linguaggio naturale al codice SQL non sia più una sfida puramente
sintattica, ma un complesso problema di gestione del contesto e del
ragionamento logico. Le recenti scoperte nell'ambito dei Large Language
Models, unite a metodologie avanzate come il Chain-of-Thought e il
Self-Correction, hanno permesso di raggiungere livelli di precisione
(Execution Accuracy) inimmaginabili fino a pochi anni fa. Nonostante i
progressi compiuti dalle soluzioni commerciali e dai recenti studi
accademici, l'analisi evidenzia il permanere di alcune criticità
strutturali. In primo luogo, emerge una forte dipendenza dal contesto:
l'efficacia di molti strumenti decade sensibilmente in presenza di
schemi di database vasti e complessi, qualora non supportati da
meccanismi di recupero (RAG) adeguatamente ottimizzati. Parallelamente,
si riscontra una carenza in termini di robustezza operativa: la capacità
dei sistemi di gestire autonomamente gli errori di runtime restituiti
dal DBMS appare spesso limitata o priva della necessaria trasparenza per
l'utente finale. Infine, un limite significativo riguarda la
flessibilità di orchestrazione. Raramente i tool attuali offrono una
personalizzazione profonda del flusso di lavoro, che permetta di
separare nettamente fasi distinte come la gestione dello schema, la
generazione della query e la successiva validazione dei risultati,
impedendo così un controllo granulare su ogni passaggio del processo.
Proprio in questo spazio si inserisce il lavoro di questa tesi. Sulla
base delle fondamenta teoriche e tecnologiche analizzate in questo
capitolo, nei capitoli successivi verrà presentata un'architettura di
Query Orchestrator progettata per massimizzare l'affidabilità della
traduzione attraverso un sistema modulare di gestione dello schema e un
ciclo iterativo di raffinamento basato sui feedback in tempo reale del
database.